\documentclass[11pt,a4paper]{article}
\usepackage[margin=2.5cm]{geometry}
\usepackage{amsmath,amssymb,amsthm,mathtools}
\usepackage{hyperref}
\usepackage{xcolor}
\usepackage{booktabs}
\usepackage{enumitem}
\usepackage{tcolorbox}
\tcbuselibrary{theorems}

\hypersetup{colorlinks=true, linkcolor=blue!60!black, citecolor=blue!60!black}

\newtheorem{theorem}{Theorem}
\newtheorem{proposition}[theorem]{Proposition}
\newtheorem{lemma}[theorem]{Lemma}
\newtheorem{corollary}[theorem]{Corollary}
\newtheorem{definition}[theorem]{Definition}
\newtheorem{remark}[theorem]{Remark}

\newtcolorbox{keyresult}{colback=blue!5!white, colframe=blue!50!black,
  fonttitle=\bfseries, title=Key Result}

\DeclareMathOperator{\proj}{proj}
\DeclareMathOperator{\tr}{tr}
\DeclareMathOperator{\diag}{diag}

\newcommand{\R}{\mathbb{R}}
\newcommand{\E}{\mathbb{E}}
\newcommand{\calL}{\mathcal{L}}
\newcommand{\calX}{\mathcal{X}}
\newcommand{\calO}{\mathcal{O}}
\newcommand{\bX}{\mathbf{X}}
\newcommand{\bx}{\mathbf{x}}

\title{\textbf{Rigorous Gradient Derivation for the $\Omega$-Framework}\\[0.5em]
\large With First-Order Taylor Decomposition and\\Comparison to Giannou et al.\ (2022)}
\author{Eugene Shcherbinin}
\date{March 2026}

\begin{document}
\maketitle

\tableofcontents
\newpage

% ══════════════════════════════════════════════════════════════════
\section{Setup and Notation}
% ══════════════════════════════════════════════════════════════════

\subsection{The Data Pipeline}

Fix the following objects from the $\Omega$-framework (Items 1--19):

\begin{center}
\begin{tabular}{cll}
\toprule
\textbf{Symbol} & \textbf{Space} & \textbf{Description} \\
\midrule
$\calX$ & $\in O$ & Observed reality (raw experience) \\
$X = \tau(\calX)$ & $\in \Pi = W \cup S$ & Proposition (rationalised observation) \\
$\bX = \phi_\psi(X)$ & $\in \R^m$ & Embedding (numerical representation) \\
$Y = f(\bX)$ & $\in \R^k$ & True causal map (unknown) \\
$\hat{Y} = \hat{f}_\theta(\bX)$ & $\in \R^k$ & Learned approximation \\
$\calL(Y, \hat{Y})$ & $\in \R$ & Loss function \\
\midrule
$\theta$ & $\in \R^l$ & Model parameters \\
$\psi$ & $\in \R^p$ & Embedding parameters \\
$\beta = (\theta, \psi)$ & $\in \R^{l+p}$ & All learnable parameters \\
\bottomrule
\end{tabular}
\end{center}

\subsection{Assumptions}

\begin{enumerate}[label=(A\arabic*)]
\item $\calL : \R^k \times \R^k \to \R$ is continuously differentiable in both arguments.
\item $\hat{f}_\theta : \R^m \to \R^k$ is continuously differentiable in $\theta$ and in its input.
\item $\phi_\psi : \Pi \to \R^m$ is continuously differentiable in $\psi$ (via UPRT, Item 9).
\item $f : \R^m \to \R^k$ is the true causal map, continuously differentiable (when accessible).
\end{enumerate}

\subsection{Dimension Summary}

For clarity, we annotate every Jacobian with its dimensions:
\begin{align*}
\frac{\partial \calL}{\partial Y} &\in \R^{1 \times k}, &
\frac{\partial \calL}{\partial \hat{Y}} &\in \R^{1 \times k}, \\
\frac{\partial f}{\partial \bX} &\in \R^{k \times m}, &
\frac{\partial \hat{f}_\theta}{\partial \bX} &\in \R^{k \times m}, \\
\frac{\partial \hat{f}_\theta}{\partial \theta} &\in \R^{k \times l}, &
\frac{\partial \phi_\psi}{\partial \psi} &\in \R^{m \times p}.
\end{align*}

\newpage
% ══════════════════════════════════════════════════════════════════
\section{The Gradient $\nabla_\beta \log \calL$}
% ══════════════════════════════════════════════════════════════════

\subsection{Composite Structure}

The loss is:
\begin{equation}
\calL = \calL\big(f(\phi_\psi(X)),\; \hat{f}_\theta(\phi_\psi(X))\big).
\label{eq:loss}
\end{equation}

Write $\bX_\psi := \phi_\psi(X)$ for the embedding. Then $Y = f(\bX_\psi)$ and $\hat{Y} = \hat{f}_\theta(\bX_\psi)$, so:
\[
\calL = \calL(Y(\psi), \hat{Y}(\theta, \psi)).
\]
Note: $Y$ depends on $\psi$ (through the embedding) but \emph{not} on $\theta$. $\hat{Y}$ depends on both.

\subsection{Gradient of $\log \calL$}

By the chain rule for $\log$:
\begin{equation}
\boxed{\nabla_\beta \log \calL = \frac{1}{\calL} \cdot \nabla_\beta \calL}
\label{eq:log-grad}
\end{equation}

So we need $\nabla_\theta \calL$ and $\nabla_\psi \calL$.

\subsection{Gradient w.r.t.\ $\theta$ (the model parameters)}

Since $Y = f(\bX_\psi)$ does not depend on $\theta$:
\begin{equation}
\boxed{
\nabla_\theta \calL = \underbrace{\frac{\partial \calL}{\partial \hat{Y}}}_{1 \times k} \cdot \underbrace{\frac{\partial \hat{f}_\theta(\bX_\psi)}{\partial \theta}}_{k \times l} \;\in\; \R^{1 \times l}
}
\label{eq:grad-theta}
\end{equation}

This is \textbf{standard backpropagation} (Item 19, eq.~1).

Therefore:
\begin{equation}
\nabla_\theta \log \calL = \frac{1}{\calL} \cdot \frac{\partial \calL}{\partial \hat{Y}} \cdot \frac{\partial \hat{f}_\theta}{\partial \theta}.
\label{eq:log-grad-theta}
\end{equation}

\subsection{Gradient w.r.t.\ $\psi$ (the embedding parameters)}

Both arguments of $\calL$ depend on $\psi$ through $\bX_\psi = \phi_\psi(X)$. By the multivariate chain rule:

\begin{equation}
\boxed{
\nabla_\psi \calL = \bigg[\underbrace{\frac{\partial \calL}{\partial Y} \cdot \frac{\partial f}{\partial \bX}}_{\text{through } f \text{ (causal)}} + \underbrace{\frac{\partial \calL}{\partial \hat{Y}} \cdot \frac{\partial \hat{f}_\theta}{\partial \bX}}_{\text{through } \hat{f}_\theta \text{ (learned)}}\bigg] \cdot \frac{\partial \phi_\psi}{\partial \psi} \;\in\; \R^{1 \times p}
}
\label{eq:grad-psi}
\end{equation}

\textbf{The two-term structure} (Item 19, eq.~2):
\begin{itemize}[nosep]
\item \textbf{Term A} (causal): $\frac{\partial \calL}{\partial Y} \cdot \frac{\partial f}{\partial \bX} \cdot \frac{\partial \phi_\psi}{\partial \psi}$ — propagates through the \emph{true} causal map $f$. \textbf{Inaccessible} if $f$ is unknown.
\item \textbf{Term B} (learned): $\frac{\partial \calL}{\partial \hat{Y}} \cdot \frac{\partial \hat{f}_\theta}{\partial \bX} \cdot \frac{\partial \phi_\psi}{\partial \psi}$ — propagates through the \emph{learned} approximation $\hat{f}_\theta$. Always computable.
\end{itemize}

Therefore:
\begin{equation}
\nabla_\psi \log \calL = \frac{1}{\calL}\bigg[\frac{\partial \calL}{\partial Y} \cdot \frac{\partial f}{\partial \bX} + \frac{\partial \calL}{\partial \hat{Y}} \cdot \frac{\partial \hat{f}_\theta}{\partial \bX}\bigg] \cdot \frac{\partial \phi_\psi}{\partial \psi}.
\label{eq:log-grad-psi}
\end{equation}

\subsection{The Full Gradient $\nabla_\beta \log \calL$}

Combining \eqref{eq:log-grad-theta} and \eqref{eq:log-grad-psi}:

\begin{keyresult}
Let $\beta = (\theta, \psi) \in \R^{l+p}$. Then:
\begin{equation}
\nabla_\beta \log \calL = \frac{1}{\calL}
\begin{pmatrix}
\dfrac{\partial \calL}{\partial \hat{Y}} \cdot \dfrac{\partial \hat{f}_\theta}{\partial \theta} \\[1.5em]
\bigg[\dfrac{\partial \calL}{\partial Y} \cdot \dfrac{\partial f}{\partial \bX} + \dfrac{\partial \calL}{\partial \hat{Y}} \cdot \dfrac{\partial \hat{f}_\theta}{\partial \bX}\bigg] \cdot \dfrac{\partial \phi_\psi}{\partial \psi}
\end{pmatrix}
\label{eq:full-gradient}
\end{equation}
The $\theta$-block has dimension $l$. The $\psi$-block has dimension $p$. The $\psi$-block contains the inaccessible term $\partial f / \partial \bX$.
\end{keyresult}


\newpage
% ══════════════════════════════════════════════════════════════════
\section{First-Order Taylor Expansion and the Trinity Decomposition}
\label{sec:taylor}
% ══════════════════════════════════════════════════════════════════

\subsection{Taylor Expansion of the Loss}

Let $Y^* \in \R^k$ be a reference prediction (e.g., the Bayes-optimal predictor at a stationary $\theta^*$). Define the residual $\delta = \hat{Y} - Y^*$. Taylor-expanding $\calL$ in its second argument:

\begin{equation}
\calL(Y, Y^* + \delta) = \calL(Y, Y^*) + \nabla_{\hat{Y}} \calL(Y, Y^*)^\top \delta + \frac{1}{2}\delta^\top \nabla^2_{\hat{Y}} \calL(Y, Y^*) \, \delta + \calO(\|\delta\|^3).
\label{eq:taylor-loss}
\end{equation}

\subsection{The Expected Risk and Trinity Decomposition}

The expected risk (Item 20):
\begin{equation}
R(\theta, \psi) = \E_{\calX \sim p}\big[\calL\big(f(\phi_\psi(\tau(\calX))),\; \hat{f}_\theta(\phi_\psi(\tau(\calX)))\big)\big].
\label{eq:risk}
\end{equation}

By the Trinity partition (Item 21), since $\tau : O \to \Pi$ is a bijection and $\Pi = W \cup S$:

\begin{equation}
\boxed{
R(\theta, \psi) = P_W \cdot \E[\calL \mid X \in W] + P_S \cdot \E[\calL \mid X \in S]
}
\label{eq:trinity-risk}
\end{equation}

where $P_W = P_\theta(\tau^{-1}(W))$ and $P_S = P_\theta(\tau^{-1}(S))$.

\subsection{Taylor + Trinity = Loss Decomposition}

Substituting \eqref{eq:taylor-loss} into \eqref{eq:trinity-risk} and writing $H_{\hat{Y}} = \nabla^2_{\hat{Y}}\calL(Y, Y^*)$:

\begin{keyresult}
\begin{align}
R(\theta, \psi) &\approx \underbrace{P_W \cdot \E[\calL(Y, Y^*) \mid W] + P_S \cdot \E[\calL(Y, Y^*) \mid S]}_{\text{irreducible loss at reference}} \nonumber \\[0.5em]
&\quad + \underbrace{P_W \cdot \E[\nabla_{\hat{Y}}\calL^\top \delta \mid W] + P_S \cdot \E[\nabla_{\hat{Y}}\calL^\top \delta \mid S]}_{\text{first-order correction (bias)}} \label{eq:trinity-taylor} \\[0.5em]
&\quad + \underbrace{\tfrac{1}{2}P_W \cdot \E[\delta^\top H_{\hat{Y}} \delta \mid W] + \tfrac{1}{2}P_S \cdot \E[\delta^\top H_{\hat{Y}} \delta \mid S]}_{\text{second-order correction (variance / curvature)}} \nonumber
\end{align}
\end{keyresult}

\textbf{Interpretation} (this is the key entry point):
\begin{enumerate}[nosep]
\item \textbf{Row 1}: Irreducible error — the loss you'd get even with the best possible predictor. Includes Pirsig's ``Quality loss'' (information destroyed by $\tau$).
\item \textbf{Row 2}: Bias — directional misalignment of $\hat{Y}$ from $Y^*$. Vanishes at a stationary point. Decomposes separately over $W$ and $S$: the learner may be biased differently on truths vs.\ falsehoods.
\item \textbf{Row 3}: Variance/curvature — the cost of dispersion around $Y^*$. The Hessian $H_{\hat{Y}}$ controls how sensitive the loss is to prediction errors. Also decomposes over $W$ and $S$.
\end{enumerate}

\subsection{Gradient of the Trinity-Expanded Risk}

Differentiating \eqref{eq:trinity-risk} w.r.t.\ $\theta$, and noting that $P_W, P_S$ may depend on $\theta$ if the sampling distribution $p_\theta$ does:

\begin{equation}
\nabla_\theta R = P_W \, \E[\calL \cdot s_\theta + \nabla_\theta \calL \mid W] + P_S \, \E[\calL \cdot s_\theta + \nabla_\theta \calL \mid S]
\label{eq:trinity-grad}
\end{equation}

where $s_\theta(\calX) := \nabla_\theta \log p_\theta(\calX)$ is the \textbf{score function}.

Each conditional expectation contains both the REINFORCE term ($\calL \cdot s_\theta$, exploration) and the backprop term ($\nabla_\theta \calL$, exploitation). The Trinity splits each into a $W$-component and an $S$-component.


\newpage
% ══════════════════════════════════════════════════════════════════
\section{Parameter-Space Taylor Expansion}
% ══════════════════════════════════════════════════════════════════

Alternatively, expand $R$ in $\theta$ around a stationary point $\theta^*$ where $\nabla_\theta R(\theta^*) = 0$. Let $\delta_\theta = \theta - \theta^*$:

\begin{equation}
R(\theta^* + \delta_\theta) \approx R(\theta^*) + \frac{1}{2}\delta_\theta^\top \Big(\underbrace{H_W^\theta}_{W\text{-curvature}} + \underbrace{H_S^\theta}_{S\text{-curvature}}\Big) \delta_\theta
\label{eq:param-taylor}
\end{equation}

where $H_A^\theta = \nabla^2_\theta R_A(\theta^*)$ for $A \in \{W, S\}$. The eigenspectra of $H_W^\theta$ and $H_S^\theta$ reveal whether the learner's parameter sensitivity comes from getting truths wrong or from getting falsehoods wrong.


\newpage
% ══════════════════════════════════════════════════════════════════
\section{Comparison with Giannou et al.\ (2022)}
\label{sec:giannou}
% ══════════════════════════════════════════════════════════════════

Giannou et al.\ study the convergence of policy gradient methods to Nash equilibria in general stochastic games. Their framework:

\begin{itemize}[nosep]
\item $N$-player stochastic game $\mathcal{G} = (\mathcal{S}, \mathcal{N}, \{\mathcal{A}_i\}, \{R_i\}, P, \zeta, \rho)$
\item Policy gradient update: $\pi_{n+1} = \proj_\Pi(\pi_n + \gamma_n \hat{v}_n)$
\item Gradient dominance property (GDP): $V_{i,\rho}(\pi'_i; \pi_{-i}) - V_{i,\rho}(\pi_i; \pi_{-i}) \leq C_\mathcal{G} \max_{\tilde\pi_i} \langle \nabla_i V_i, \tilde\pi_i - \pi_i \rangle$
\item SOS condition: $(\pi - \pi^*)^\top \text{Jac}_v(\pi^*)(\pi - \pi^*) < 0$
\end{itemize}

\subsection{Where the $\Omega$-Framework is Superior}

\begin{center}
\renewcommand{\arraystretch}{1.4}
\begin{tabular}{p{3.5cm}p{5cm}p{5cm}}
\toprule
\textbf{Dimension} & \textbf{Giannou et al.} & \textbf{$\Omega$-Framework} \\
\midrule
\textbf{Gradient structure} & Single gradient $v(\pi) = \nabla_i V_i$. & Five-component gradient (eq.\ 27 of theos): exploration + exploitation + evidence-seeking + alignment + opponent-shaping. \\

\textbf{Representation} & Fixed tabular parameterization $\pi_i : \mathcal{S} \to \Delta(\mathcal{A}_i)$. & Learnable embedding $\phi_\psi$ with its own gradient (eq.~\ref{eq:grad-psi}). Creates a bilevel optimization that Giannou doesn't address. \\

\textbf{Loss decomposition} & Single value function $V_{i,\rho}(\pi)$. & Trinity decomposition: $R = P_W \E[\calL \mid W] + P_S \E[\calL \mid S]$. Convergence can be tracked \emph{separately} on truths vs.\ falsehoods. \\

\textbf{GDP refinement} & Single GDP inequality (their Lemma 2). & Trinity-refined GDP with separate $W$ and $S$ bounds. Rates may differ: the learner converges faster on truths than falsehoods (or vice versa). \\

\textbf{Irreducible error} & Not modeled — assumes the game is fully specified. & Explicit irreducible loss from the $O$-set (Pirsig's Quality, G\"odel's incompleteness). The gap $Y - \hat{Y}$ has a floor that no gradient can eliminate. \\

\textbf{Multi-agent structure} & Agents are symmetric (same ``type'' of gradient). & G\"odelian complementarity: agents with different $F_i$ have structurally different blind spots $B_i$. Multi-agent $>$ single-agent not by speed but by \emph{coverage}. \\

\textbf{Evidence quality} & Not modeled. & Keynesian weight of evidence $V_i(\calX)$ as second loss axis. Evidence-seeking gradient term (eq.\ 11 of theos): $-\mu V^{-2} \nabla_\theta V$. \\

\textbf{Convergence type} & Local convergence to Nash with probability $1 - \delta$ (their Thm 1). & Same guarantee applies to the exploitation term. But the exploration term (REINFORCE) enables \emph{global} search over $O$. \\
\bottomrule
\end{tabular}
\end{center}

\subsection{Formal Refinement: Trinity GDP}

\begin{proposition}[Trinity-Refined Gradient Dominance]
\label{prop:trinity-gdp}
Under the $\Omega$-framework, for any policy profile $\pi = (\pi_i)_{i \in \mathcal{N}}$ and unilateral deviation $\pi'_i$:
\begin{equation}
V_{i,\rho}(\pi'_i; \pi_{-i}) - V_{i,\rho}(\pi_i; \pi_{-i}) \leq C_\mathcal{G}\Big[P_W \max_{\tilde\pi_i}\langle \nabla_i V_i^W(\pi), \tilde\pi_i - \pi_i\rangle + P_S \max_{\tilde\pi_i}\langle \nabla_i V_i^S(\pi), \tilde\pi_i - \pi_i\rangle\Big]
\label{eq:trinity-gdp}
\end{equation}
where $V_i^W, V_i^S$ are the value functions restricted to states reachable through $W$- and $S$-propositions respectively.
\end{proposition}

\begin{proof}[Proof sketch]
Apply Giannou's Lemma 2 separately to the $W$ and $S$ components of the reward, using the Trinity decomposition of the expected trajectory reward:
\[
V_{i,s}(\pi) = \E\bigg[\sum_{t=0}^{T(\tau)} R_i(s_t, \alpha_t) \bigg| s_0 = s\bigg] = \underbrace{\E\big[R_i \cdot \mathbf{1}_{X_t \in W}\big]}_{V_i^W} + \underbrace{\E\big[R_i \cdot \mathbf{1}_{X_t \in S}\big]}_{V_i^S}.
\]
The GDP (which follows from the linearity of $V$ in $\pi_i$) applies to each component since $V_i^W + V_i^S = V_i$.
\end{proof}

\begin{corollary}[Asymmetric Convergence Rates]
If $\pi^*$ is SOS and the Trinity-refined GDP holds, then the convergence rate from Giannou's Theorem 2 applies separately to $V_i^W$ and $V_i^S$, potentially with different constants $\mu_W, \mu_S$ in the SOS condition. In particular, a learner may:
\begin{itemize}[nosep]
\item Converge at rate $\calO(1/n)$ on $W$-propositions (truths it gets right),
\item Converge at rate $\calO(1/\sqrt{n})$ on $S$-propositions (falsehoods it gets wrong),
\end{itemize}
and the overall rate is dominated by the slower component.
\end{corollary}


\subsection{The Five-Component Gradient vs.\ Giannou's PG Template}

Giannou's PG template (their eq.\ PG):
\[
\pi_{n+1} = \proj_\Pi(\pi_n + \gamma_n \hat{v}_n)
\]
with $\hat{v}_n = v(\pi_n) + U_n + b_n$ (true gradient + noise + bias).

In the $\Omega$-framework (Item 52), $\hat{v}_n$ decomposes as:
\begin{equation}
\hat{v}_{i,n} = \underbrace{\omega_i \E[\calL \cdot s_{\theta_i}]}_{\text{(a) exploration}} + \underbrace{\omega_i \E[\nabla_{\theta_i} \calL]}_{\text{(b) exploitation}} + \underbrace{\omega_i \mu(-V_i^{-2}\nabla_{\theta_i} V_i)}_{\text{(c) evidence}} + \underbrace{\lambda \nabla_{\theta_i} D}_{\text{(d) alignment}} + \underbrace{\sum_{j \neq i} \frac{\partial R_i}{\partial \theta_j}\frac{\partial \theta'_j}{\partial \theta_i}}_{\text{(e) opponent shaping}}
\label{eq:five-component}
\end{equation}

\textbf{What Giannou captures}: Only term (b), the exploitation gradient. Their $v(\pi_n)$ is the exact policy gradient, which is the exploitation term.

\textbf{What Giannou misses}:
\begin{itemize}[nosep]
\item \textbf{(a) Exploration}: The REINFORCE score-function term that adjusts \emph{where} the agent looks. This is not a noise term ($U_n$) — it's a systematic gradient component.
\item \textbf{(c) Evidence}: The Keynesian evidence-seeking term. Absent from any standard MARL framework.
\item \textbf{(d) Alignment}: The consensus penalty. Related to communication in Dec-POMDPs but not in Giannou's model.
\item \textbf{(e) Opponent shaping}: The LOLA term. Giannou explicitly assumes agents treat others as stationary (their Model 1). The $\Omega$-framework models opponent adaptation.
\end{itemize}

\begin{remark}
Giannou's convergence results (Theorems 1--3) apply to the exploitation term (b) in isolation. The $\Omega$-framework's contribution is showing that in practice, agents should optimize \emph{all five terms simultaneously}, and the convergence guarantees from Giannou give a lower bound on what the exploitation term alone achieves.
\end{remark}


\newpage
% ══════════════════════════════════════════════════════════════════
\section{The SOS Condition and Trinity Curvature}
% ══════════════════════════════════════════════════════════════════

Giannou's SOS condition (their eq.\ SOS):
\[
(\pi - \pi^*)^\top \text{Jac}_v(\pi^*)(\pi - \pi^*) < 0 \quad \text{for all } \pi \in \Pi \setminus \{\pi^*\}.
\]

In the $\Omega$-framework, the Jacobian $\text{Jac}_v$ decomposes via the Trinity:

\begin{equation}
\text{Jac}_v(\pi^*) = P_W \cdot \text{Jac}_{v^W}(\pi^*) + P_S \cdot \text{Jac}_{v^S}(\pi^*).
\label{eq:trinity-jac}
\end{equation}

This means the SOS condition can be \emph{refined}:

\begin{definition}[Trinity-SOS]
A Nash policy $\pi^*$ is \textbf{Trinity-SOS} if:
\begin{align}
(\pi - \pi^*)^\top \text{Jac}_{v^W}(\pi^*)(\pi - \pi^*) &< 0 \quad \text{(SOS on truths)} \label{eq:sos-W} \\
(\pi - \pi^*)^\top \text{Jac}_{v^S}(\pi^*)(\pi - \pi^*) &< 0 \quad \text{(SOS on falsehoods)} \label{eq:sos-S}
\end{align}
for all $\pi \neq \pi^*$ sufficiently close.
\end{definition}

\begin{remark}
Trinity-SOS is \emph{stronger} than Giannou's SOS (it implies SOS by linearity) but gives \emph{finer} convergence information. The eigenvalues of $\text{Jac}_{v^W}$ and $\text{Jac}_{v^S}$ independently control convergence on truths vs.\ falsehoods. A policy might be SOS overall but unstable on $S$ (the falsehood component has a positive eigenvalue), meaning the learner is systematically worsening on falsehoods while improving on truths — a phenomenon Giannou's framework cannot detect.
\end{remark}


\newpage
% ══════════════════════════════════════════════════════════════════
\section{Summary: What the $\Omega$-Framework Adds}
% ══════════════════════════════════════════════════════════════════

\begin{enumerate}
\item \textbf{The gradient is richer}: Five components vs.\ one. Giannou's convergence results give a floor for what exploitation alone achieves; the $\Omega$-framework shows the full gradient agents should compute.

\item \textbf{The loss decomposes}: The Trinity $W/S$ split gives separate convergence tracking on truths vs.\ falsehoods. This is invisible in Giannou's framework.

\item \textbf{There's an irreducible floor}: The $O$-set (G\"odel, Pirsig, Keynes) means $R(\theta^*) > 0$ for any $\theta^*$. Giannou's framework implicitly assumes the game is fully specified; the $\Omega$-framework models what is \emph{structurally unknowable}.

\item \textbf{Representation learning matters}: The $\psi$-gradient (eq.~\ref{eq:grad-psi}) with its inaccessible term $\partial f / \partial \bX$ is absent from Giannou's tabular setting. This is where Pearl's causal structure (Item 46--48) provides the missing $\partial f / \partial \bX$ via DAG-structured backpropagation.

\item \textbf{Multi-agent is structurally different}: G\"odelian complementarity means multi-agent systems can know strictly more than any single agent — not by being faster (Giannou's concern) but by having non-overlapping blind spots. The five-component gradient (eq.~\ref{eq:five-component}) captures this via the opponent-shaping term.

\item \textbf{The Taylor expansion is the entry method}: Equation \eqref{eq:trinity-taylor} is the key decomposition. At any stationary point, the risk separates into irreducible loss (row 1), bias (row 2), and variance (row 3), each further split over $W$ and $S$. This is the diagnostic tool: which component of the loss, on which partition of reality, is the learner failing?
\end{enumerate}


\newpage
% ══════════════════════════════════════════════════════════════════
\section*{Appendix: Lean 4 Formalization Summary}
% ══════════════════════════════════════════════════════════════════

A Lean 4 file \texttt{Theos.lean} accompanies this document, formalizing:
\begin{itemize}[nosep]
\item The truth value type \texttt{TruthVal} and partition $W, O, S$ (Items 2--4)
\item The Trinity exhaustiveness and disjointness theorems (proved)
\item $O = \Pi^c$ (proved)
\item The bijection $\tau : O \to \Pi$ (axiomatized)
\item The UPRT: $\Pi \leftrightarrow \R^\infty$ (axiomatized)
\item The learning pipeline: $\texttt{LossFunction}$, $\texttt{Embedding}$, $\texttt{CausalMap}$, $\texttt{Approximator}$, $\texttt{LearningPipeline}$ (defined as structures)
\item The gradient decompositions (axiomatized as structural claims)
\item Connection to Giannou: the $\Omega$-gradient subsumes Giannou's PG template (stated)
\end{itemize}

The file requires Lean 4 with Mathlib. Proofs marked \texttt{sorry} require analysis-level Mathlib (continuous differentiability, Fr\'echet derivatives) which is available but verbose to formalize fully.

\end{document}
